\vspace{-5pt}
\section{引言}
鸟瞰视图（BEV）识别模型~\cite{OFT,pseudo-lidar,VPN,LSS,bevformer,CVT,PETR} 在自动驾驶领域引起了广泛关注，因为它们能够将来自多传感器的局部原始观测自然地整合到统一的三维输出空间中。
典型的BEV模型建立在图像骨干网络（Image Backbone）之上，后接一个视图变换（View Transformation）模块，将透视图特征（Perspective Feature）提升为BEV特征，再经由BEV特征编码器（Encoder）和任务特定头（Task-specific Head）进一步处理。
尽管大量工作致力于设计视图变换模块~\cite{CVT,LSS,bevformer} 并将不断增长的下游任务~\cite{LSS,fiery} 纳入新的识别框架中，BEV模型中图像骨干网络的研究却鲜受关注。
作为一个前沿且需求极高的领域，将现代图像骨干网络引入自动驾驶是自然而然的方向。
然而，研究社区仍坚持使用 VoVNet~\cite{Vovnet} 以享受其大规模深度预训练（Depth Pre-training）~\cite{DD3D} 带来的优势。
本文专注于释放现代图像特征提取器（Image Feature Extractor）在BEV识别中的全部潜力，为未来研究者在该领域探索更好的图像骨干网络设计打开大门。

然而，仅使用这些现代图像骨干网络而不进行适当的预训练无法产生满意的结果。
例如，在 ImageNet \cite{ImageNet} 上预训练的 ConvNeXt-XL \cite{ConvNext} 骨干网络在3D目标检测（3D Object Detection）上的表现仅与在 DDAD-15M 上预训练的 VoVNet-99 \cite{DD3D} 相当，尽管后者参数量是前者的3.5倍。
本文将现代图像骨干网络适配困难的原因归结为以下问题：
1) 自然图像与自动驾驶场景之间的领域差异。
在通用2D识别任务上预训练的骨干网络难以感知3D场景，尤其是深度估计（Depth Prediction）。
2) 当前BEV检测器复杂的结构。
以 BEVFormer~\cite{bevformer} 为例。3D边界框和物体类别标签的监督信号被视图编码器和物体解码器（Decoder）与图像骨干网络隔离开来，而每个编码器和解码器均由多层 Transformer 组成。
用于将通用2D图像骨干网络适配到自动驾驶任务的梯度流被堆叠的 Transformer 层扭曲。

为克服上述在将现代图像骨干网络适配到BEV识别中的困难，本文将透视监督（Perspective Supervision）引入 BEVFormer，即来自透视图任务的额外监督信号直接作用于骨干网络。
透视监督引导骨干网络学习2D识别任务中缺失的3D知识，并克服BEV检测器的复杂性，显著促进模型优化。
具体而言，本文在骨干网络上构建了一个透视3D检测头（Perspective 3D Detection Head）\cite{DD3D}，其以图像特征为输入并直接预测目标物体的3D边界框和类别标签。
该透视图头的损失（称为透视损失（Perspective Loss））作为辅助检测损失（Auxiliary Detection Loss），与来自BEV头的原始损失（BEV损失）相加。
两个检测头及其对应的损失项被联合训练。
此外，本文自然地将两个检测头结合成一个两阶段（Two-stage）BEV检测器——\textbf{BEVFormer v2}。
由于透视图头功能完整，能够在透视图上生成高质量的物体提议（Object Proposal），本文将其作为第一阶段提议。
将这些提议编码为物体查询（Object Query），与原始 BEVFormer 中的可学习查询合并，形成混合物体查询（Hybrid Object Query），再送入第二阶段检测头生成最终预测。

本文开展了广泛的实验来验证所提透视监督的有效性和必要性。
透视损失有助于图像骨干网络的适配，从而提升检测性能并加速模型收敛。
而若无此监督，即使采用更长的训练计划，模型也无法达到可比较的结果。
最终，本文成功将现代图像骨干网络适配到BEV模型中，在 nuScenes \cite{nuscenes} 测试集上达到63.4\% NDS。

本文贡献总结如下：
 \begin{itemize}
     \item 指出透视监督是将通用2D图像骨干网络适配到BEV模型的关键，通过透视图上的检测损失显式地添加此监督。
     \item 提出一种新颖的两阶段BEV检测器 BEVFormer v2，包含一个透视3D检测头和一个BEV检测头，前者的提议与后者的物体查询相结合。
     \item 通过与最新图像骨干网络结合，在 nuScenes 数据集上取得了相比先前SOTA结果的显著提升，验证了方法的有效性。
 \end{itemize}
